%%
%% This is file `example-inline.tex',
%% generated with the docstrip utility.
%%
%% The original source files were:
%%
%% exsol.dtx  (with options: `example-inline')
%% 
%% This is a generated file.
%% 
%% Copyright (C) 2019 by Walter Daems <walter.daems@uantwerpen.org>
%% 
%% This file may be distributed and/or modified under the conditions of
%% the LaTeX Project Public License, either version 1.3 of this license
%% or (at your option) any later version.  The latest version of this
%% license is in:
%% 
%%    http://www.latex-project.org/lppl.txt
%% 
%% and version 1.3 or later is part of all distributions of LaTeX version
%% 2005/12/01 or later.
%% 
\documentclass[11pt,a4paper]{article}

\usepackage[german]{babel}
\usepackage[inline,usesolutionserieslabels]{exsol}
\usepackage{enumitem}

\setlength{\exsolexercisetopbottomsep}{0pt plus 0pt minus 1pt}
\setlength{\exsolexerciseleftmargin}{2em}
\setlength{\exsolexerciserightmargin}{1em}
\setlength{\exsolexerciseparindent}{0em}
\setlength{\exsolexerciselabelsep}{1ex}
\setlength{\exsolexerciselabelwidth}{30pt}
\setlength{\exsolexerciseitemindent}{0pt}
\setlength{\exsolexerciseparsep}{\parskip}

\title{Local inline example, from the \textsf{ExSol} package}
\author{Philippe Marti}
\setlength{\parindent}{0em}
\begin{document}
\maketitle

\section{Gleichungssysteme und Geraden}

Ein bisschen Theorie\ldots

\begin{exerciseseries}[solsubrule=\hrule]{Gleichungssysteme}
  \begin{exercise}
    Die Summe zweier Zahlen ist 17 und ihre Differenz 7. Bestimme die
    beiden Zahlen!
  \end{exercise}
  \begin{solution}
    5 und 12
  \end{solution}

  \begin{exercise}
    Die Differenz einer Zahl und dem Dreifachen einer zweiten Zahl ist
    14. Bestimme die beiden Zahlen, falls die zweite Zahl ein Zehntel
    der ersten ist.
  \end{exercise}
  \begin{solution}
    20 und 2
  \end{solution}
\end{exerciseseries}
~\\
Etwas mehr Theorie\ldots

\begin{exerciseseries}{Geraden}
  \begin{exercise}
    Berechne den Schnittpunkt von \mbox{$y=3x+1$} und \mbox{$y=3x-7$}.
  \end{exercise}
  \begin{solution}
    Es gibt keinen Schnittpunkt
  \end{solution}

  \begin{exercise}
    Die Familie Meier fordert Offerten f\"ur eine Heizungsreparatur
    ein. Firma A berechnet f\"ur die Fahrtkosten Fr. 42.- und f\"ur
    jede Arbeitsstunde 76.-. Bei der Firma B sind die Fahrtkosten
    Fr. 35.- und jede Arbeitsstunde wird mit Fr. 80.- berechnet.
    \begin{enumerate}[label=\alph*)]
    \item Welche Kosten entstehen f\"ur beide Firmen, wenn ein Monteur
      3.5 Stunden
      f\"ur die Arbeit benötigt? Welche Firma ist in diesem Fall
      kostengünstiger?

    \item Wie lauten die Gleichungen derjenigen zwei linearen
      Funktionen, die jeder Arbeitszeit $x$ (in Stunden) die
      entstehenden Kosten $y$ (in Franken) zuordnet?

    \item Berechne, bei welcher Arbeitszeit die Kosten bei beiden
      Firmen gleich sind.
    \end{enumerate}
  \end{exercise}
  \begin{solution}
    \begin{enumerate}[label=\alph*)]
    \item Firma A: 308.- $|$ Firma B: 315.-
    \item A: $y=76x+42$ $|$ B: $y=80x+35$
    \item Bei $1\frac{3}{4}$ Stunden
    \end{enumerate}
  \end{solution}
\end{exerciseseries}

\end{document}
\endinput
%%
%% End of file `example-inline.tex'.
